\documentclass[10pt]{article}

\usepackage[utf8]{inputenc}

\usepackage[T1]{fontenc}

\usepackage{amsmath}

\usepackage{amsfonts}

\usepackage{amssymb}

\usepackage[version=4]{mhchem}

\usepackage{stmaryrd}

\usepackage{bbold}

\usepackage{graphicx}

\usepackage[export]{adjustbox}

\graphicspath{ {./images/} }



\begin{document}

Лum.: [1] А х и е з е р Н. И., Г л а зм а н И. М., Теория गинейных операторов в гильбертовом пространстве, 3 изд., Јекции по функциональному анализу, пер: с франц., 2 изд., М.,

1979.

В. К. Соболев.



ӘРМИТОВА МАТРИЦА, ә р м и т о в о - с и м ме трическая мат риц а, с а мосо п р я ж нная м а т р и ц а,- квадратная матрица $A=\left\|a_{i k}\right\|$ над полем $\mathbb{C}$, совпадаюшшая со своей э р м и т в ос о п я ж н но й м а т и ц е й



\[

A^{*}=\bar{A}^{\top}=\left\|{\overline{a_{k}}}\right\|

\]



т. е. матрица, әлементы к-рой удовлетворяют условию $a_{i k}=\bar{a}_{k i}$. Если все $a_{i k} \in \mathbb{R}$, то Э. м. $A-$-имметрич. матрица. Э. м. фиксированного порядка образуют векторное пространство над полем $\mathbb{R}$. Если $A, B-$ Э. м. одного порядка, то $A B+B A-\vartheta$. м. Относительно операции $A \circ B=1 / 2(A B+B A)$ Э. м. (порядка $n$ ) coставляют йорданову алаебру. Произведение $A B$ Э. м. является Э.м. тогда и только тогда, когда $A$ и $B$ перестановочны.



Э. м. порядка $n$ служат матрицами әрмитовых преобразований $n$-мерного унитарного пространства в ортонормированном базисе (см. Самосопряженное линейное преобразование ). С другой стороны, Э. м.- это матрицы әрмитовых форм в $n$-мерном комплексном векторном пространстве. Как и әрмитовы боржы, Э. м. могут быть определены над любым телом c антиинволюцией.



Bсе собственные значения Ә. м. действительны. Д ля каждой Э. м. $A$ существует унитарная матрица $U$ такая, что $U^{-1} A U$ - диагональная действительная матрица. Ә. м. наз. н е о т р и ц а т е лы ь о й (или



\includegraphics[max width=\textwidth, center]{2023_07_09_2399ba1e9fcc3a0b960eg-1(3)}

если все ее главные миноры неотрицательны, и п о л ожи т е льн о п р е д е л е н н о й, если все они положительны. Неотрицательные (положительно определенные) Ә. м. соответствуют нөотрицатөльным (положительно определенным) эрмитовым линейным преобразованиям и әрмитовым формам. А. Л. оницик.



ЭРМИТОВА МЕТРИКА - 1) Э. м. В к о м пл е кс ном в е т о р ном прост р а н т в е ложительно определенная әрмитова форла в $V$. Пространство $V$, снабженное Э. м., наз, у н и т а р ны м (или к о м п л е к с но е в к ли До в ы м, или ә рмитовым векторным) прост ранство м, а Ә. м. в немІІ р о и з е д е н и е м. Любые две Э. м. в $V$ переводятся друг в друга автоморфизмом пространства $V$. Таким образом, множество всех Э. м. в $V$ является однородным пространством группы $G L_{n}(\mathbb{C})$ и отождествляется $G L_{n}(\mathbb{C}) / U(n), \quad$ где $n=\operatorname{dim} V$.



Комплексное векторное пространство $V$ можно рассматривать как вещественное векторное пространство $V^{\mathbb{R}}$, снабженное оператором комплексной структуры



\includegraphics[max width=\textwidth, center]{2023_07_09_2399ba1e9fcc3a0b960eg-1}

ется евклидовой метрикой (скалярным произведением) в $V$, а форма $\omega=\operatorname{Im} h$ - невырожденной кососимметрической билинейной формой в $V$. При этом $g(J x, J y)=$ $=g(x, y), \quad \omega(J x, J y)=\omega(x, y), \omega(x, y)=g(x, J y)$. JI юбая из форм $g, \omega$ однозначно определяет $h$.



\begin{enumerate}

  \setcounter{enumi}{1}

  \item Э. м. в ком пл лксно к ве к то рн о м

\end{enumerate}



\includegraphics[max width=\textwidth, center]{2023_07_09_2399ba1e9fcc3a0b960eg-1(2)}

базе $M$, сопоставляющая точке $p \in M$ әрмитову метрику $g_{p}$ в слое $E(p)=\pi^{-1}(p)$ расслоения л и удовлетворяюшая следуюшему условию гладкости: для любых гладіих Јокальных сечений $e, e^{\prime}$ расслоения л функция $p \longmapsto g_{p}\left(e_{p}, e_{p}^{\prime}\right)$ является гладкой.



В Јюбом комплексном векторном расслоөнии существует Э. м. Связность $\nabla$ в комплексном векторном расслоении $\pi$ наз. согласованної с $Э$. м. оператор $J$ комплексної структуры в слоях расслоения $\pi$ ковариантно постоянны (т. с. $\nabla g=\nabla J=0$ ), иначе говоря, если соответствующий параллельный перенөс слоев расслоения $\pi$ вдоль кривых на базе является изоморфизмом слоев как унитарных пространств. Для любой Ә. м. существует согласованная с ней связность, к-рая, вообще говоря, не единственна. В случае, когда $\pi$ есть голоморфоное векторное расслоение над комплексным многообразием $M$ (см. Векторное аналитическое расслоение), существует единственная связность $\nabla$ расслоения $\pi$, согласованная с данной Ә. М. и удовлетворяющая следующему условию: ковариантная производная любого голоморфоного сечения $e$ расслоения $\pi$ относительно любого антиголоморфного комплексного векторного поля $\bar{X}$ на $M$ равна нулю



\includegraphics[max width=\textwidth, center]{2023_07_09_2399ba1e9fcc3a0b960eg-1(1)}

Форму кривизны этой связности можно рассматривать как 2-форму типа $(1,1)$ на $M$ со значением в расслоении эндоморфизмов векторного расслоения л. Канонич. связность можно рассматривать также как связность в главном $G L(n, \mathbb{C})$-расслоении $\tilde{\pi}: P \rightarrow M$, ассоциированном с голоморфным расслоөнием $\pi$ комплексной размерности $n$. Она характеризуется как единственная связность в $\tilde{\pi}$, горизонтальные подпространства к-рой являются комплексными подпространствами касательных пространств комшлексного многообразия $P$.



Лum.: [1] К о б а я с и ШІ,, Н о м и д з у К., Основы дифференциальной геометрии, пер.' с англ., т. 2, М., 1981; [2] Л и хн е р о в и ч А., Теория связностей в целом и группы голономий, пер. с франц., М., 1960; [3] у ә л л с Р., Диффференциальное исчисление на комплексных многообразиях, пер. с англ., М., $1976 ;$ [4] В е й л ь А., Введение в теорию кәлеровых многообразий, пер. с франц., М., 1961.

М. В. Алексеевский.



ӘРМИТОВА СВЯЗНОСТБ - аффинная связность на әрмитовом многообразии $M$, относительно к-рой ковариантно постоянны тензор $\varphi$ комплексной структуры и фундаментальная 2-форма $\Omega=\frac{i}{2} g_{\alpha \beta}{ }^{\beta} \wedge \overline{\omega^{\alpha}}$, следовательно и әрмитова форма $d s^{2}=g_{\alpha} \bar{\omega}^{\alpha} \omega^{\beta}$. Если афффинная связность на $M$ задана локальными формами связности $\omega_{\beta}^{\alpha}=\Gamma_{\beta}^{\alpha} \omega^{\gamma}+\Gamma_{\beta}^{\alpha} \overline{\omega^{\gamma}}$, то эти условия выражаются в виде $\omega_{\beta}^{\frac{\alpha}{\beta}}=\omega_{\beta}^{\bar{\alpha}}=0, \omega_{\bar{\beta}}^{\bar{\alpha}}=\bar{\omega}_{\beta}^{\alpha}, d g_{\alpha \beta}=$ $=\bar{\omega}_{\alpha}^{\gamma} g_{\gamma \beta}+g_{\alpha} \omega_{\beta}^{\gamma}$. На заданном әрмитовом многообразии $M$ существует одна и только одна Э. с., у к-рой $\Gamma_{\beta \bar{\gamma}}^{\alpha}=0$



Обобщением Ә. с. является почт и ә р и то в а с в я зность, к-рая определяется аналогичными условиями ковариантного постоянства тензоров $\varphi_{j}^{i}$ и $g_{i j}$ с $g_{k l} \varphi_{i}^{k} \varphi_{j}^{l}=g_{i j}$ на почти эрмитовом многообразии $\tilde{M}$. Почти Э. с. на заданном $\tilde{M}$ существует с нек-рым произволом, к-рый сосредоточен в ее тензоре кручения: если тензоры кручения двух почти Э. с. совпадают, то и сами связности совпадают. Например, существует одна и только одна почти Э.с., у к-рой формы кручения являются суммами «чистых» форм (т. е. форм тига $(2,0)$ и типа $(0,2)),-$ вторая каноническая связность Лихнеровича.



Лит.: [1] Л и х н е р о в и ч А., Теория связностей в целом и группы голономий, пер. с франц., М., 1960; [2] Y a n o K., Differential geometry on complex and almost complex spaces, Oxf., 1965.

Ю. $\Gamma$. Jулисте.



ЭРМИТОВА СТРУКТУРА на многообразии $M$ - пара $(J, g)$, состояшая пз комплексной структуры $J$ многообразия $M$ и әрмитовой метрики $g$ в касательном расслоении $T M$, т. е. римановой метрики $g$, инвариантной относительно $J$ :



\[

g(J X, J Y)=g(X, Y)

\]



для любых векторных полей $X, Y$ на $M$. Э. с. задает в каждом касательном пространстве $T_{p} M$ структуру эрмитова векторного пространства (см. Эрлитова метрика). Многообразие с Ә. с. наз. э р м и т о в ы м м н о г о о б р а з и м. Э. с. определяет на $M$ дифффе-





\end{document}